\section{Free Will and the Laws of Physics}

On a blog by a soi-disant reactionary, we read:

\begin{quotex}
As for free will, a minute's thought will show that the notion of free will implies an abrogation of the laws of physics. 
\flright{\textsc{Dennis Mangan}, 2009-08-01, blog is now defunct}
\end{quotex}

Such an ``abrogation'' sounds pretty serious, so I thought about it a minute, pondering which particular law of physics would be abrogated. First I thought of \textbf{Newton}, whose theories were considered to be ``laws of physics'' for more than a couple of hundred years. Unfortunately, his ``laws'' were definitively refuted by Einstein, so Newton's theory turned out not to be a law, but is instead a brilliantly conceived, albeit false, recreation of physics.

So, perhaps he means the \textbf{Theory of Relativity}. Can that theory predict whether I'll choose pistachio or strawberry ice cream next Sunday? Well, according to that theory, every body follows a ``world line'' in space-time, so presumably a physicist could compute my world line and predict which flavour I will choose. But another minute's thought will convince you that what we have is not a ``law of physics'', but rather a hypothesis that demands empirical testing. So far, at any rate, no one has computed my world-line, indeed, no one has even proposed how that could possibly be done. Hence, until there is an experiment that can be tested, it is premature to consider it a law of physics. Furthermore, another minute's reflection will open up this question: am ``I'' truly a material body subject to laws of physics? Yes, I certainly have a material body, but does that exhaust who or what I am? It appears that the laws of physics simply assume I am nothing but a material object, rather than prove it. Therefore, it cannot be the Theory of Relativity.

Perhaps, then, he means \textbf{Quantum Theory}, which, for example, can predict the distribution of electrons passing through a slit, although it cannot predict the course of any particular electron. Presumably, then, quantum physics can predict what percentage of pistachio or strawberry ice cream I'll choose over the course of the year, but not what I'll choose on any given Sunday.

But I can do a similar thing. For example, I can produce a very accurate chart of the probability distribution of the next million throws of the dice at a casino, although I cannot tell you what the next throw will be with any certainty. Yet, alas, I can do that without invoking any laws of physics at all. In particular, I don't need to know anything about the mechanics of the motion of the dice---it is simply a matter of mathematics, not physics.

So more than a minute's reflection opens up the following questions:

\begin{enumerate}


\item Why should we suppose there are ``laws of physics''?

\item Does physics account for the totality of reality?

\item Is a scientific law anything other than a clever creation of the human mind?
\end{enumerate}



\flrightit{Posted on 2009-09-30 by Cologero}

